\documentclass[12pt,a4paper]{article}

% ---------------------------------------------------------
% Pacotes
% ---------------------------------------------------------
\usepackage[utf8]{inputenc}
\usepackage[T1]{fontenc}
\usepackage[portuguese]{babel}
\usepackage{geometry}
\geometry{margin=2.5cm}
\usepackage{graphicx}
\usepackage{hyperref}
\usepackage{listings}
\usepackage{xcolor}
\usepackage{booktabs}
\usepackage{enumitem}
\usepackage{amsmath}
\usepackage{float}
\usepackage{caption}
\usepackage{subcaption}
\usepackage{fancyhdr}
\usepackage{titlesec}
\usepackage{parskip}
\usepackage{microtype}

% ---------------------------------------------------------
% Cores e estilo de código
% ---------------------------------------------------------
\definecolor{codegray}{rgb}{0.95,0.95,0.95}
\definecolor{codeblue}{rgb}{0.13,0.29,0.53}
\definecolor{codegreen}{rgb}{0.0,0.5,0.0}
\definecolor{codepurple}{rgb}{0.55,0.0,0.55}
\definecolor{codered}{rgb}{0.6,0.0,0.0}

\lstdefinestyle{pythonstyle}{
    backgroundcolor=\color{codegray},
    basicstyle=\ttfamily\footnotesize,
    keywordstyle=\color{codeblue}\bfseries,
    commentstyle=\color{codegreen}\itshape,
    stringstyle=\color{codered},
    numberstyle=\tiny\color{gray},
    numbers=left,
    stepnumber=1,
    breaklines=true,
    frame=single,
    rulecolor=\color{lightgray},
    tabsize=4,
    showspaces=false,
    showstringspaces=false,
    language=Python,
    morekeywords={dataclass, Optional, Literal, list, dict, set, tuple},
}

\lstset{style=pythonstyle}

% ---------------------------------------------------------
% Cabeçalho e rodapé
% ---------------------------------------------------------
\pagestyle{fancy}
\fancyhf{}
\rhead{PRI --- Projeto 1}
\lhead{Motor de Pesquisa sobre o Repositorium}
\cfoot{\thepage}
\renewcommand{\headrulewidth}{0.4pt}

% ---------------------------------------------------------
% Formatação de títulos de secções
% ---------------------------------------------------------
\titleformat{\section}{\large\bfseries}{}{0em}{\thesection\quad}
\titleformat{\subsection}{\normalsize\bfseries}{}{0em}{\thesubsection\quad}

% ---------------------------------------------------------
% Metadados do documento
% ---------------------------------------------------------
\title{
    \vspace{-1cm}
    {\large Universidade do Minho}\\[0.3cm]
    {\large Mestrado em Engenharia Informática}\\[0.3cm]
    \rule{\linewidth}{0.4pt}\\[0.5cm]
    {\LARGE \textbf{Processamento e Recuperação de Informação}}\\[0.3cm]
    {\Large Projeto 1 --- Motor de Pesquisa sobre o Repositorium}\\[0.3cm]
    \rule{\linewidth}{0.4pt}
}
\author{
    Daniela Rodrigues \and Rodrigo [Apelido]\\[0.2cm]
    \small Universidade do Minho
}
\date{Maio de 2026}

% ==========================================================
\begin{document}
% ==========================================================

\maketitle
\thispagestyle{empty}
\newpage

\tableofcontents
\newpage

% ==========================================================
\section{Introdução}
% ==========================================================

Este relatório descreve o desenvolvimento de um motor de pesquisa de informação (IR) aplicado ao \textbf{Repositorium}, o repositório institucional da Universidade do Minho. O projeto foi desenvolvido no âmbito da unidade curricular de Processamento e Recuperação de Informação (PRI) e tem como objetivo construir, de raiz, um sistema capaz de indexar, pré-processar e pesquisar publicações científicas.

O sistema é composto por vários módulos interligados:
\begin{itemize}
    \item um \textbf{scraper} baseado em Selenium para recolha automática de metadados;
    \item uma \textbf{base de dados SQLite} para persistência dos documentos e do índice;
    \item um \textbf{pipeline de pré-processamento NLP} com suporte a stemming e lematização;
    \item um \textbf{índice invertido} com skip pointers para pesquisa booleana eficiente;
    \item um \textbf{motor TF-IDF} com implementações própria e via scikit-learn;
    \item um \textbf{classificador Naïve Bayes} para categorização automática de documentos;
    \item uma \textbf{API REST} desenvolvida em FastAPI com frontend HTML.
\end{itemize}

A escolha do Repositorium prende-se com a sua relevância no contexto académico português e com a disponibilidade pública dos seus metadados, o que facilita a avaliação da qualidade dos resultados de pesquisa.

% ==========================================================
\section{Recolha de Dados --- Scraper}
% ==========================================================

\subsection{Motivação e abordagem}

O Repositorium é uma aplicação Angular (DSpace 8), o que significa que o conteúdo é renderizado dinamicamente pelo browser --- uma simples chamada HTTP não devolve o HTML útil. Para contornar esta limitação, optámos por usar o \textbf{Selenium WebDriver} com o Chrome em modo \emph{headless}, que permite aguardar pela renderização do Angular antes de extrair os elementos.

\subsection{Arquitetura do scraper}

A classe principal \texttt{UMinhoDSpace8Scraper} encapsula toda a lógica de navegação. O processo decorre em três fases:

\begin{enumerate}
    \item \textbf{Descoberta de links}: a função \texttt{collect\_all\_links()} percorre a paginação da coleção, extraindo os URLs de cada publicação até atingir o limite configurado (\texttt{MAX\_ITEMS}).
    \item \textbf{Extração de metadados}: para cada URL, a função \texttt{get\_paper\_info()} navega até à página \texttt{/full} do item e lê a tabela de metadados Dublin Core (\texttt{dc.title}, \texttt{dc.date.issued}, \texttt{dc.contributor.author}, etc.).
    \item \textbf{Persistência}: os resultados são guardados em \texttt{scraper\_results.json}.
\end{enumerate}

\begin{table}[H]
\centering
\caption{Campos extraídos pelo scraper para cada publicação}
\begin{tabular}{@{}lll@{}}
\toprule
\textbf{Campo Dublin Core} & \textbf{Chave JSON} & \textbf{Descrição} \\
\midrule
\texttt{dc.title}               & \texttt{title}         & Título da publicação \\
\texttt{dc.date.issued}         & \texttt{year}          & Ano de publicação \\
\texttt{dc.identifier.doi}      & \texttt{doi}           & Identificador DOI \\
\texttt{dc.contributor.author}  & \texttt{authors}       & Lista de autores \\
\texttt{dc.description.abstract}& \texttt{abstract}      & Resumo \\
---                             & \texttt{document\_link}& Link para o PDF \\
\bottomrule
\end{tabular}
\end{table}

\subsection{Configuração de ambiente}

O scraper inclui uma função \texttt{find\_chrome\_executable()} que deteta automaticamente a instalação do Chrome em Windows, Linux e macOS, eliminando a necessidade de configuração manual. Caso o Chrome não esteja instalado, é lançada uma exceção clara com instruções.

% ==========================================================
\section{Armazenamento --- Base de Dados SQLite}
% ==========================================================

\subsection{Schema relacional}

A base de dados foi desenhada para suportar os requisitos de armazenamento do projeto, seguindo um modelo relacional normalizado. O schema inclui as seguintes tabelas principais:

\begin{itemize}
    \item \textbf{\texttt{documents}}: armazena o documento com o seu conteúdo bruto (\texttt{raw\_content}) e o texto pré-processado serializado em JSON (\texttt{processed\_text}).
    \item \textbf{\texttt{authors}}: tabela normalizada de autores, com índice único no nome completo para evitar duplicados.
    \item \textbf{\texttt{document\_authors}}: tabela de junção N:N entre documentos e autores, com campo de posição para preservar a ordem dos autores.
    \item \textbf{\texttt{document\_metadata}}: pares chave/valor flexíveis para metadados adicionais (sujeito, tipo de documento, etc.).
    \item \textbf{\texttt{document\_versions}}: histórico de versões do conteúdo de cada documento.
    \item \textbf{\texttt{index\_terms}} e \textbf{\texttt{index\_postings}}: representação persistente do índice invertido.
\end{itemize}

\subsection{Classe DocumentStore}

A classe \texttt{DocumentStore} centraliza todas as operações de leitura e escrita. Utiliza um \emph{context manager} para gestão de conexões e suporta transações explícitas. As operações mais relevantes são:

\begin{itemize}
    \item \texttt{save\_documents(docs, processed\_texts)}: insere documentos em lote, resolve autores por upsert e persiste metadados.
    \item \texttt{save\_inverted\_index(index)}: serializa o índice invertido na base de dados.
    \item \texttt{get\_stats()}: devolve contagens por tabela para diagnóstico.
\end{itemize}

% ==========================================================
\section{Pré-processamento de Texto}
% ==========================================================

\subsection{Pipeline NLP}

O módulo \texttt{preprocessor.py} implementa um pipeline de pré-processamento configurável, construído sobre o NLTK. A classe \texttt{Preprocessor} aceita uma instância de \texttt{PreprocessorConfig} que controla cada etapa individualmente.

O pipeline executa as seguintes etapas, por esta ordem:

\begin{enumerate}
    \item \textbf{Tokenização}: usando \texttt{word\_tokenize} do NLTK, com suporte a inglês e português.
    \item \textbf{Normalização de caixa}: conversão para minúsculas.
    \item \textbf{Remoção de pontuação}: filtro de tokens sem caracteres alfanuméricos.
    \item \textbf{Remoção de stopwords}: usando os léxicos NLTK para inglês e português em simultâneo.
    \item \textbf{Lematização}: usando \texttt{WordNetLemmatizer} com etiquetagem POS (para inglês); para português, lematização simples como nomes.
    \item \textbf{Stemming}: Porter Stemmer (inglês) ou Snowball Stemmer (português).
    \item \textbf{Filtro de comprimento mínimo}: remoção de tokens com menos de 2 caracteres.
\end{enumerate}

\subsection{Configurações disponíveis}

O sistema disponibiliza três \emph{factory functions} para os casos de uso mais comuns:

\begin{lstlisting}[caption={Factory functions para configurações pré-definidas}]
def make_stemming_preprocessor(language="english") -> Preprocessor:
    # Stemming sem lematizacao (mais rapido, menor vocabulario)
    return Preprocessor(PreprocessorConfig(
        use_stemming=True, use_lemmatisation=False, remove_stopwords=True
    ))

def make_lemmatisation_preprocessor(language="english") -> Preprocessor:
    # Lematizacao sem stemming (mais precisa, vocabulario maior)
    return Preprocessor(PreprocessorConfig(
        use_stemming=False, use_lemmatisation=True, remove_stopwords=True
    ))

def make_bare_preprocessor(language="english") -> Preprocessor:
    # Apenas tokenizacao e lowercase, sem filtragem
    return Preprocessor(PreprocessorConfig(
        use_stemming=False, use_lemmatisation=False, remove_stopwords=False
    ))
\end{lstlisting}

\subsection{Comparação stemming vs.\ lematização}

A principal diferença entre as duas abordagens está na forma como os termos são reduzidos à sua forma base:

\begin{table}[H]
\centering
\caption{Comparação entre stemming e lematização}
\begin{tabular}{@{}lll@{}}
\toprule
\textbf{Aspeto}       & \textbf{Stemming}          & \textbf{Lematização} \\
\midrule
Velocidade           & Rápido                     & Mais lento (requer POS) \\
Vocabulário          & Menor                      & Maior \\
Qualidade dos termos & Radicais (podem não existir) & Formas válidas \\
Exemplo: ``running'' & \texttt{run}               & \texttt{run} \\
Exemplo: ``better''  & \texttt{better}            & \texttt{good} \\
\bottomrule
\end{tabular}
\end{table}

% ==========================================================
\section{Índice Invertido}
% ==========================================================

\subsection{Estrutura de dados}

O índice invertido é implementado no módulo \texttt{inverted\_index.py} e assenta em duas estruturas principais:

\begin{itemize}
    \item \textbf{\texttt{Posting}}: regista, para um par (termo, documento), a frequência do termo (\texttt{tf}) e as posições de ocorrência no texto (lista de inteiros). As posições são usadas em pesquisas de proximidade.
    \item \textbf{\texttt{PostingsList}}: agrega todos os postings de um termo, a sua frequência de documento (\texttt{df}) e um dicionário de skip pointers.
\end{itemize}

\subsection{Skip pointers}

Para acelerar a interseção de listas de postings em queries AND, implementámos skip pointers. O intervalo de salto é calculado como $\lfloor\sqrt{n}\rfloor$, onde $n$ é o comprimento da lista:

\begin{lstlisting}[caption={Cálculo dos skip pointers}]
def build_skip_pointers(self) -> None:
    n = len(self.postings)
    if n < 3:
        self.skips = {}
        return
    skip_interval = max(1, int(math.sqrt(n)))
    self.skips = {}
    for i in range(0, n - skip_interval, skip_interval):
        self.skips[i] = i + skip_interval
\end{lstlisting}

O algoritmo de interseção com skip pointers evita comparações desnecessárias ao saltar blocos de elementos quando o doc\_id atual é menor do que o alvo:

\begin{equation}
    \text{skip\_interval} = \max\!\left(1,\ \left\lfloor\sqrt{n}\right\rfloor\right)
\end{equation}

No pior caso este algoritmo tem complexidade $O(n)$; com skip pointers, na prática aproxima-se de $O(\sqrt{n})$ em listas uniformemente distribuídas.

\subsection{Suporte a atualizações incrementais}

A função \texttt{add\_document()} permite indexar um documento sem reconstruir o índice completo, e atualiza os skip pointers apenas nas listas de postings afetadas.

% ==========================================================
\section{TF-IDF e Similaridade}
% ==========================================================

\subsection{Modelo de pesquisa vetorial}

O módulo \texttt{tfidf.py} implementa um motor de pesquisa baseado no modelo vetorial. Cada documento e cada query são representados como vetores no espaço dos termos, ponderados por TF-IDF. A relevância é medida pela similaridade de cosseno entre os vetores.

\subsection{Cálculo TF-IDF}

O sistema suporta três esquemas de TF:

\begin{table}[H]
\centering
\caption{Esquemas de TF suportados}
\begin{tabular}{@{}lll@{}}
\toprule
\textbf{Esquema} & \textbf{Fórmula} & \textbf{Quando usar} \\
\midrule
\texttt{raw}     & $\mathrm{tf}(t,d)$ & Coleções pequenas e uniformes \\
\texttt{log}     & $1 + \log(\mathrm{tf}(t,d))$ & Uso geral (padrão) \\
\texttt{boolean} & $1$ se $\mathrm{tf}>0$, senão $0$ & Pesquisa booleana simples \\
\bottomrule
\end{tabular}
\end{table}

O IDF é calculado com suavização para evitar divisão por zero:

\begin{equation}
    \mathrm{idf}(t) = \log\!\left(\frac{N+1}{\mathrm{df}(t)+1}\right) + 1
\end{equation}

onde $N$ é o número total de documentos e $\mathrm{df}(t)$ é o número de documentos que contêm o termo $t$.

A similaridade de cosseno entre dois vetores $\mathbf{v}_1$ e $\mathbf{v}_2$ é:

\begin{equation}
    \cos(\mathbf{v}_1, \mathbf{v}_2) = \frac{\mathbf{v}_1 \cdot \mathbf{v}_2}{\|\mathbf{v}_1\| \cdot \|\mathbf{v}_2\|}
\end{equation}

\subsection{Dois backends}

O \texttt{TFIDFEngine} suporta dois backends que produzem resultados comparáveis:

\begin{itemize}
    \item \textbf{Custom}: implementação própria com vetores esparsos em dicionários Python. Mais transparente e auditável.
    \item \textbf{Sklearn}: usa \texttt{TfidfVectorizer} e matrizes \texttt{scipy.sparse}. Mais eficiente em coleções grandes.
\end{itemize}

Ambos os backends são construídos em \texttt{load} e podem ser alternados em tempo de execução via parâmetro da API.

\subsection{Matriz de similaridade}

O método \texttt{similarity\_matrix()} calcula a matriz $N \times N$ de similaridades par-a-par entre todos os documentos indexados. O método \texttt{similar\_to(doc\_id)} usa esta lógica para devolver os $K$ documentos mais próximos de um dado documento.

% ==========================================================
\section{Pesquisa Booleana}
% ==========================================================

\subsection{Sintaxe suportada}

O módulo \texttt{boolean\_search.py} implementa um motor de pesquisa booleana completo. Suporta os operadores \texttt{AND}, \texttt{OR}, \texttt{NOT} com a precedência correta ($\text{NOT} > \text{AND} > \text{OR}$), agrupamento por parênteses, e \texttt{AND} implícito entre termos separados por espaço.

Exemplos de queries válidas:
\begin{itemize}
    \item \texttt{"information retrieval"} → \texttt{information AND retrieval}
    \item \texttt{"neural OR deep AND learning"} → \texttt{neural OR (deep AND learning)}
    \item \texttt{"machine NOT learning"} → documentos com ``machine'' mas sem ``learning''
    \item \texttt{"(neural OR deep) AND learning"} → agrupamento explícito
\end{itemize}

\subsection{Parser de queries}

O módulo \texttt{query\_utils.py} implementa um parser de descida recursiva que constrói uma \textbf{Abstract Syntax Tree} (AST) da query. Os nós da AST (\texttt{QueryNode}) têm um operador (\texttt{AND}/\texttt{OR}/\texttt{NOT}/\texttt{TERM}) e filhos esquerdo/direito.

A avaliação da AST percorre a árvore recursivamente, realizando as interseções, uniões e diferenças das listas de postings correspondentes.

Adicionalmente, o módulo implementa \textbf{expansão de query} via WordNet: para cada termo da query, são adicionados sinónimos (configurável, máximo de 2 por defeito), alargando o recall da pesquisa.

% ==========================================================
\section{Classificação de Documentos --- Naïve Bayes}
% ==========================================================

\subsection{Abordagem}

O módulo \texttt{classifier.py} implementa um classificador Multinomial Naïve Bayes para categorização automática de publicações em cinco áreas temáticas: \emph{Computer Science}, \emph{Health/Medicine}, \emph{Engineering}, \emph{Social Sciences} e \emph{Mathematics}.

\subsection{Estratégia de treino semi-supervisionada}

Como o dataset do Repositorium não vem com etiquetas de categoria, adotámos uma estratégia de \emph{bootstrapping}:

\begin{enumerate}
    \item \textbf{Atribuição de etiquetas iniciais}: um dicionário de palavras-chave por categoria (\texttt{CATEGORIES}) é usado para etiquetar automaticamente os documentos. Documentos que não acumulem matches suficientes ficam na categoria ``other''.
    \item \textbf{Treino supervisionado}: o \texttt{MultinomialNB} do scikit-learn é treinado sobre as etiquetas geradas, usando \texttt{TfidfVectorizer} sobre os tokens pré-processados.
    \item \textbf{Predição}: novos documentos são classificados pelo modelo treinado, com \emph{fallback} para as palavras-chave em caso de confiança baixa.
\end{enumerate}

\subsection{Avaliação}

O classificador calcula as métricas standard de classificação: \emph{precision}, \emph{recall} e F1-score por categoria. Opcionalmente, realiza validação cruzada com 5 \emph{folds} (controlado pelo parâmetro \texttt{CLASSIFIER\_RUN\_CROSS\_VALIDATION} na configuração).

\begin{table}[H]
\centering
\caption{Categorias e exemplos de palavras-chave associadas}
\begin{tabular}{@{}ll@{}}
\toprule
\textbf{Categoria}    & \textbf{Exemplos de palavras-chave} \\
\midrule
computer\_science  & algorithm, machine learning, neural, database, nlp \\
health\_medicine   & clinical, patient, cancer, genomics, epidemiology \\
engineering        & mechanical, sensor, renewable, embedded, circuit \\
social\_sciences   & education, policy, economics, psychology, ethics \\
mathematics        & probability, algebra, differential, bayesian, theorem \\
\bottomrule
\end{tabular}
\end{table}

% ==========================================================
\section{API REST e Frontend}
% ==========================================================

\subsection{FastAPI}

A API é implementada em \texttt{fastapi\_app.py} usando o framework \textbf{FastAPI}. Os endpoints principais são:

\begin{table}[H]
\centering
\caption{Endpoints da API REST}
\begin{tabular}{@{}lll@{}}
\toprule
\textbf{Método} & \textbf{Endpoint}       & \textbf{Descrição} \\
\midrule
GET  & \texttt{/search}            & Pesquisa TF-IDF; aceita \texttt{q}, \texttt{top\_k}, \texttt{backend} \\
GET  & \texttt{/boolean}           & Pesquisa booleana com \texttt{AND}/\texttt{OR}/\texttt{NOT} \\
GET  & \texttt{/stats}             & Estatísticas do índice \\
GET  & \texttt{/documents/\{id\}}  & Metadados de um documento \\
GET  & \texttt{/similar/\{id\}}    & Documentos mais similares \\
GET  & \texttt{/classify}          & Classifica um documento dado o seu ID \\
\bottomrule
\end{tabular}
\end{table}

O engine é carregado uma vez no arranque da aplicação e mantido em memória. Um mecanismo de \emph{lazy reload} (\texttt{ensure\_loaded()}) permite recarregar os dados se o caminho do ficheiro de dados mudar --- útil em contexto de testes.

\subsection{Frontend}

O frontend é uma single-page application em HTML/CSS/JavaScript puro, servida pela própria API em \texttt{/ui}. Apresenta uma caixa de pesquisa, opções para escolher o backend (custom vs.\ sklearn) e o tipo de pesquisa (TF-IDF vs.\ booleana), e uma lista de resultados com \emph{snippets} do abstract.

% ==========================================================
\section{Avaliação de Desempenho}
% ==========================================================

O módulo \texttt{evaluation.py} implementa um framework de avaliação de desempenho para o sistema:

\begin{itemize}
    \item \textbf{Tempo de indexação}: mede o tempo de construção do índice invertido e do TF-IDF para stemming vs.\ lematização.
    \item \textbf{Uso de memória}: usa \texttt{tracemalloc} para medir o pico de uso de memória durante a indexação.
    \item \textbf{Tempo de resposta a queries}: mede a latência de cada pesquisa.
    \item \textbf{Relevância dos resultados}: calcula Precision@K e MAP usando um conjunto de \emph{relevance judgements} (\emph{qrels}) definido manualmente.
\end{itemize}

As \emph{qrels} incluem queries de referência como ``information retrieval'', ``machine learning'' e ``deep learning'', com títulos de documentos relevantes definidos por substring match.

% ==========================================================
\section{Configuração Centralizada}
% ==========================================================

Todos os parâmetros configuráveis do sistema estão centralizados em \texttt{config.py}, num \emph{dataclass} \texttt{Settings} que age como \emph{singleton}. Os valores podem ser sobrescritos por variáveis de ambiente com o prefixo \texttt{IR\_}:

\begin{lstlisting}[caption={Exemplo de override por variável de ambiente}]
# Porta da API
export IR_API_PORT=9000

# Esquema TF
export IR_TF_SCHEME=raw

# Backend TF-IDF
export IR_DEFAULT_TFIDF_BACKEND=sklearn
\end{lstlisting}

Esta abordagem permite configurar o sistema sem tocar no código, o que facilita o deployment em diferentes ambientes (desenvolvimento, testes, produção).

% ==========================================================
\section{Testes}
% ==========================================================

O projeto inclui uma suite de testes abrangente em \texttt{pytest}, com cobertura para todos os módulos principais:

\begin{table}[H]
\centering
\caption{Ficheiros de teste e módulos cobertos}
\begin{tabular}{@{}ll@{}}
\toprule
\textbf{Ficheiro de teste}       & \textbf{Módulo testado} \\
\midrule
\texttt{test\_preprocessor.py}   & Pipeline NLP \\
\texttt{test\_index\_search.py}  & Índice invertido e pesquisa \\
\texttt{test\_tfidf.py}          & Motor TF-IDF \\
\texttt{test\_boolean\_search.py} & Pesquisa booleana \\  % via test_query_utils
\texttt{test\_query\_utils.py}   & Parser e expansão de query \\
\texttt{test\_classifier.py}     & Classificador Naïve Bayes \\
\texttt{test\_database.py}       & Camada de persistência \\
\texttt{test\_scraper.py}        & Scraper (com mocks) \\
\texttt{test\_api.py}            & Endpoints FastAPI \\
\bottomrule
\end{tabular}
\end{table}

Os testes do scraper usam \emph{mocking} do WebDriver para não depender de uma ligação à internet, e os testes da API usam o \texttt{TestClient} do FastAPI com um dataset de teste reduzido.

% ==========================================================
\section{Estrutura do Repositório}
% ==========================================================

\begin{verbatim}
PRI_PROJETO1/
├── src/
│   ├── config.py               # Configuração centralizada
│   ├── scraper/
│   │   ├── scraper.py          # Selenium scraper
│   │   ├── main.py             # Entrypoint do scraper
│   │   └── scraper_results.json
│   ├── storage/
│   │   ├── database.py         # Camada SQLite
│   │   └── populate_db.py      # Script de população
│   ├── search/
│   │   ├── preprocessor.py     # Pipeline NLP
│   │   ├── inverted_index.py   # Índice invertido
│   │   ├── tfidf.py            # Motor TF-IDF
│   │   ├── boolean_search.py   # Motor booleano
│   │   ├── query_utils.py      # Parser e expansão
│   │   ├── classifier.py       # Classificador NB
│   │   ├── evaluation.py       # Avaliação
│   │   └── nlp.py              # Utilitários NLP
│   ├── api/
│   │   └── fastapi_app.py      # API REST
│   └── frontend/
│       └── index.html          # Interface web
├── data/
│   └── repositorium.db         # Base de dados SQLite
├── tests/                      # Suite de testes pytest
├── requirements.txt
└── conftest.py
\end{verbatim}

% ==========================================================
\section{Conclusões}
% ==========================================================

O sistema desenvolvido cobre os requisitos do projeto de forma completa, implementando desde a recolha de dados até à interface de pesquisa. As principais contribuições são:

\begin{itemize}
    \item Um scraper robusto para o DSpace 8 com Angular, capaz de navegar paginação e extrair metadados estruturados.
    \item Um pipeline NLP configurável com suporte a inglês e português, comparando stemming e lematização.
    \item Um índice invertido com skip pointers que acelera as interseções de postings em queries booleanas complexas.
    \item Um motor TF-IDF com dois backends (custom e sklearn) que podem ser comparados diretamente.
    \item Um classificador automático baseado em Naïve Bayes com estratégia de bootstrapping para datasets sem anotação prévia.
    \item Uma API REST bem estruturada com frontend integrado.
\end{itemize}

Como trabalho futuro, seria interessante explorar modelos de linguagem mais recentes (e.g., embeddings BERT) para melhorar a qualidade do ranking, bem como implementar técnicas de \emph{query expansion} mais sofisticadas baseadas em modelos de linguagem.

% ==========================================================
% Referências
% ==========================================================
\newpage
\begin{thebibliography}{9}

\bibitem{manning2008}
Manning, C.D., Raghavan, P., \& Schütze, H. (2008).
\textit{Introduction to Information Retrieval}.
Cambridge University Press.
\url{https://nlp.stanford.edu/IR-book/}

\bibitem{nltk2009}
Bird, S., Klein, E., \& Loper, E. (2009).
\textit{Natural Language Processing with Python}.
O'Reilly Media.

\bibitem{sklearn}
Pedregosa et al. (2011).
Scikit-learn: Machine Learning in Python.
\textit{Journal of Machine Learning Research}, 12, 2825--2830.

\bibitem{fastapi}
Ramírez, S. (2019).
FastAPI. \url{https://fastapi.tiangolo.com/}

\bibitem{repositorium}
Universidade do Minho.
Repositorium --- Repositório Institucional da Universidade do Minho.
\url{https://repositorium.sdum.uminho.pt/}

\bibitem{dspace}
DSpace Community (2023).
\textit{DSpace 8 Documentation}.
\url{https://dspace-labs.github.io/DSpace/}

\end{thebibliography}

\end{document}
